\documentclass[12pt,a4paper]{article}
\usepackage[utf8]{inputenc}
\usepackage[turkish]{babel}
\usepackage{amsmath}
\usepackage{graphicx}
\usepackage{hyperref}
\usepackage{listings}
\usepackage{xcolor}
\usepackage{geometry}
\geometry{a4paper, margin=1in}

% Define colors for code listings
\definecolor{codegreen}{rgb}{0,0.6,0}
\definecolor{codegray}{rgb}{0.5,0.5,0.5}
\definecolor{codepurple}{rgb}{0.58,0,0.82}
\definecolor{backcolour}{rgb}{0.95,0.95,0.92}

\lstdefinestyle{mystyle}{
    backgroundcolor=\color{backcolour},   
    commentstyle=\color{codegreen},
    keywordstyle=\color{magenta},
    numberstyle=\tiny\color{codegray},
    stringstyle=\color{codepurple},
    basicstyle=\ttfamily\footnotesize,
    breakatwhitespace=false,         
    breaklines=true,                 
    captionpos=b,                    
    keepspaces=true,                 
    numbers=left,                    
    numbersep=5pt,                  
    showspaces=false,                
    showstringspaces=false,
    showtabs=false,                  
    tabsize=2
}
\lstset{style=mystyle}

\usepackage{enumitem}

\begin{document}
\shorthandoff{=}

\begin{titlepage}
    \centering
    \vspace*{1cm}
    
    {\Large\textbf{ANKARA UNIVERSITESI}}\\[0.3cm]
    {\large Muhendislik Fakultesi}\\[0.3cm]
    {\large Bilgisayar Muhendisligi Bolumu}\\[1cm]

    \begin{figure}[h]
    \centering
    \vspace{3pt}
    \includegraphics[width=5.0cm]{logo.png}
    \end{figure}


    \begin{tabular}{ll}
        \textbf{Ogrenci:} & Ahmet Akgun \\
        \textbf{Ogrenci No:} & 22290602 \\
        \textbf{Danisman:} & Enver Bagci \\
    \end{tabular}\\[1.5cm]
    
    {\Large \textbf{BLM4522 Ag Tabanli Paralel Dagitim Sistemleri}}\\[0.5cm]
    {\large \textbf{Proje Adi:} Veritabani Guvenligi ve Erisim Kontrolu}
    
    \vfill
    
    \begin{itemize}[label=\textbullet]
        \item \textbf{GitHub Depo Linki:} \url{https://github.com/aakgunnn}
        \item \textbf{Haftalık Video Takip Bağlantısı:} \url{https://www.youtube.com/@ahmetakgun359}
    \end{itemize}

\end{titlepage}
\tableofcontents
\newpage

\section{Giriş}
Bu proje kapsamında; bir veritabanı sisteminin karşılaşabileceği güvenlik tehditlerine karşı alınması gereken temel savunma önlemleri PostgreSQL ortamında incelenmiş ve simüle edilmiştir. Proje beş temel aşamadan oluşmaktadır:
\begin{enumerate}
    \item Veritabanı Şeması ve Test Verilerinin Hazırlanması
    \item Erişim Yönetimi ve Rol Bazlı Yetkilendirme (Least Privilege)
    \item Veri Şifreleme (TDE Alternatifi - pgcrypto)
    \item Denetim Kayıtları (Audit Logs)
    \item SQL Injection Zafiyeti ve Korunma Yöntemleri
\end{enumerate}

\section{Aşama 1: Ortamın Hazırlanması}
Proje testlerinin gerçekleştirilebilmesi için \texttt{security\_demo\_db} adında bir veritabanı oluşturulmuştur. Bu veritabanı içerisinde sanal bir şirketi temsil eden \texttt{company} şeması ve bu şemanın altında üç adet tablo tasarlanmıştır:
\begin{itemize}
    \item \textbf{employees:} Çalışan bilgileri (isim, departman, maaş).
    \item \textbf{users:} Sisteme giriş yapan kullanıcı bilgileri ve hashlenmiş parolalar.
    \item \textbf{financial\_records:} Hassas veri simülasyonu için kredi kartı bilgilerini barındıran tablo.
\end{itemize}

\section{Aşama 2: Erişim Yönetimi (Access Management)}
Veritabanı güvenliğinde en önemli prensiplerden biri olan ``En Az Yetki (Least Privilege)'' prensibi gereğince kısıtlı yetkilere sahip roller oluşturulmuştur.
\begin{lstlisting}[language=SQL, caption=Rol ve Kullanıcı Oluşturma]
CREATE ROLE readonly_role;
CREATE ROLE admin_role;

CREATE USER app_user WITH PASSWORD 'AppUser123!' LOGIN;
GRANT readonly_role TO app_user;

GRANT USAGE ON SCHEMA company TO readonly_role;
GRANT SELECT ON ALL TABLES IN SCHEMA company TO readonly_role;
\end{lstlisting}
Bu yapı sayesinde, uygulamanın kullandığı \texttt{app\_user} hesabı sadece veri okuma (SELECT) yetkisine sahip kılınmış, manipülasyon (INSERT, UPDATE, DELETE, DROP) işlemleri varsayılan olarak engellenmiştir.

\section{Aşama 3: Veri Şifreleme}
Hassas verilerin doğrudan düz metin (plain-text) olarak veritabanında tutulması büyük bir güvenlik riskidir. Bu projede, yapısal şifreleme için PostgreSQL'in \texttt{pgcrypto} eklentisi kullanılmış ve \texttt{financial\_records} tablosundaki \texttt{credit\_card\_info} sütunu ikili (bytes) formata çevrilerek şifrelenmiştir.
\begin{lstlisting}[language=SQL, caption=pgcrypto ile Veri Şifreleme]
CREATE EXTENSION IF NOT EXISTS pgcrypto;

ALTER TABLE company.financial_records
    ALTER COLUMN credit_card_info TYPE bytea 
    USING pgp_sym_encrypt(credit_card_info, 'SuperGizliAnahtar_123');
\end{lstlisting}
Bu işlem ile veriler disk üzerinde (Data-at-Rest) korunmuş ve sadece yetkili (anahtarı bilen) kişilerin veriyi \texttt{pgp\_sym\_decrypt} fonksiyonu üzerinden deşifre ederek okuyabilmesi sağlanmıştır.

\section{Aşama 4: Audit Logları}
Sistemde kimin, ne zaman, hangi sorguyu çalıştırdığının takibi yasal uyumluluklar ve güvenlik ihlallerinin tespiti için kritiktir. Projede detaylı audit logları aktifleştirilmiştir.
\begin{lstlisting}[language=SQL, caption=Audit Konfigürasyonu]
ALTER SYSTEM SET log_statement = 'all';
ALTER SYSTEM SET log_connections = 'on';
ALTER SYSTEM SET log_disconnections = 'on';
\end{lstlisting}
Bu sayede yetkisiz erişim denemeleri ve değiştirilen/silinen verilerin logları merkezi sunucu kayıtlarına işlenmiştir.

\section{Aşama 5: SQL Injection Zafiyeti ve Önlemler}
Web uygulamalarında sıkça karşılaşılan SQL Enjeksiyonu (SQLi) zafiyetini simüle etmek için Python (\texttt{psycopg2}) ile bir uygulama geliştirilmiştir. 

\subsection{Zafiyetli (Insecure) Durum}
Sorgular doğrudan string manipülasyonu ile birleştirildiğinde saldırgan dışarıdan mantıksal operatörler ekleyebilir:
\begin{lstlisting}[language=Python]
# Zafiyetli Kullanim
query = f"SELECT user_id, username FROM company.users WHERE username = '{username}';"
\end{lstlisting}
Saldırgan sisteme \texttt{ahmet\_it' OR '1'='1} şeklinde bir girdi gönderdiğinde, WHERE filtresi baypas edilmiş ve tablodaki tüm kullanıcıların listelenmesi sağlanmıştır.

\subsection{Korumalı (Secure) Durum - Prepared Statements}
Aynı sorgu, veri tabanına parametrik sorgular (Prepared Statements) gönderilerek revize edilmiştir.
\begin{lstlisting}[language=Python]
# Guvenli Kullanim
query = "SELECT user_id, username FROM company.users WHERE username = %s;"
cur.execute(query, (username,))
\end{lstlisting}
Bu kullanımda parametreler doğrudan argüman olarak sunucuya aktarılmış, böylece saldırganın verdiği SQL komutları işlenilebilir bir kod olarak algılanmamıştır.

\subsection{Hasar Kontrolü (Least Privilege Testi)}
Son olarak, yetki kısıtlamasının önemi test edilmiştir. Saldırgan SQL enjeksiyonu ile \texttt{DROP TABLE company.users;} komutunu çalıştırdığında; Aşama 2'de yapılandırılan \texttt{app\_user} kullanıcısının sadece SELECT yetkisi olduğu için süreç veritabanı tarafından kesilmiş ve sistem korunmuştur (Permission denied).

\section{Sonuç}
Bu projede, güvenli bir PostgreSQL veritabanı ortamının nasıl hazırlanacağı pratik şekilde uygulanmıştır. Rol tabanlı erişim kontrolü, veri şifreleme, aktivite izleme ve zafiyet engelleme (Parametrik sorgular) test edilerek doğrulanmıştır. Alınan bu önlemler, sistem mimarisinin dayanıklılığını maksimum seviyeye çıkartmaktadır.

\end{document}
